\section{背景与预备知识}\label{sec:background}
\subsection{Release--Acquire 执行关系}
一次共享内存访问表示为一个带身份的事件，记录所属线程、地址及读写类型。写 $\W(x,v)$ 向地址 $x$ 写值 $v$，读 $\R(x,v)$ 观测值 $v$。本文采用固定事件的 RA 片段：全部非初始化写为 release 写，全部读为 acquire 读；地址和写值已经确定。片段不含 RMW、fence、非原子访问或动态事件激活。每个地址有一个初始化写，所有非初始化事件组成若干线程序列。$\SB$（sequenced-before）表示各线程的严格程序顺序；实现可存其邻接边。

读自关系 $\rf$（reads-from）为每个读指定唯一的同地址写来源，并要求返回值满足该读的观测。两个写即使写入相同值，也具有不同身份，可以成为不同来源。修改序 $\co$（coherence order）在每个地址的所有写上给出严格全序，初始化写最小；它不比较不同地址的写。给定这些关系，就得到一个候选完成的执行图。

在所选片段中，读自边建立 release/acquire 同步，因果关系为
\begin{equation}\label{eq:hb}
 \hb=(\SB\cup\rf\cup\init)^+.
\end{equation}
$\init$ 将每个初始化写连到全部非初始化事件。上标 $+$ 表示非空传递闭包，因此 $\hb(e,f)$ 表示存在一条由程序顺序、同步和初始化边组成的非空路径。稀疏实现只需把初始化连到各线程首事件，再沿线程边传播。

读源和修改序还确定读取相对于更晚写的顺序。若 $r$ 读取 $w$，且 $w<_{\co}v$，则 $r$ 位于 $v$ 之前的 from-read 关系中。用关系复合 $;$ 表示先沿左关系再沿右关系，有
\begin{equation}\label{eq:eco}
 \fr=\rf^{-1};\co,\qquad
 \eco=(\rf\cup\co\cup\fr)^+.
\end{equation}
$\eco$ 汇合了同地址的读源、写序和读后写关系。本文的 RA 合法性条件为
\begin{equation}\label{eq:ra}
 \irr(\hb;\eco^?).
\end{equation}
上标 $?$ 表示加上恒等关系，$\irr$ 要求不存在从某个事件回到自身的关系对。因此式~\eqref{eq:ra} 同时排除 $\hb$ 自环，以及一条因果路径接上一条同地址一致性路径所形成的闭环。这是本文随后证明和实现共同使用的语义；RA 一致性检验中相应的读／写一致性条件可见~\cite{tunc2023}。

一个具体冲突说明这条公理如何过滤来源。若 $w\mathrel\rf r$、$w<_{\co}v$ 且 $v\mathrel\hb r$，则 $r\mathrel\fr v$，与 $v\mathrel\hb r$ 构成闭环。也就是说，读不能从一个已经被其因果前序写覆盖的旧写取值。同线程有序读取的来源也不能沿修改序倒退。这样的 coherence 约束让共同因果关系限制来源选择，具体推导在协议证明中展开。

\subsection{图分解与活跃接口}
处理候选图时，我们用一列事件集合 $B_0,B_1,\ldots$，称为\emph{袋}。它们覆盖所有事件，每条需要处理的边的两个端点至少在一个袋中共存，并且一个事件出现的袋下标形成连续区间。这是图分解的覆盖与连通性要求在线性袋序列上的形式~\cite[§3.1]{conrado2026}，称为路径分解。我们用 $K=\max_i|B_i|$ 表示最大事件袋大小；传统路径宽度为 $K-1$。编译器在处理边时保留袋中的事件关系，事件最后一次离开袋后，通过摘要保留其对未来的影响。

接口大小取决于同时活跃的端点数。一个旧写即使与很远以后的读相连，也只需要持续保留这一个端点；许多彼此不同的旧端点同时连接未来才会增大袋。对严格偏序 $L$，其覆盖关系定义为
\[
 \Cov(L)=\{(u,v)\in L:\nexists z.\ L(u,z)\land L(z,v)\}.
\]
有限偏序满足 $\Cov(L)^+=L$。这允许后文用稀疏边生成已经确定的修改序，而不把传递闭包的每一对都变成必须共存的端点。

\subsection{条件查询与知识编译}
设布尔函数 $F(X)$ 表示一组可行选择。对部分赋值 $\alpha$ 条件化，就是研究 $F\land\alpha$ 的模型；撤销一项假设则恢复相应选择。一个满足赋值可以回答“一种可行选择是什么”，却不能直接回答哪些变量取值在\emph{某个}模型中可行、哪些在\emph{所有}模型中必然成立。知识编译把构造费用前置，将 $F$ 转为支持所需查询的显式表示~\cite{darwiche2002}。

若表示引入辅助变量 $Y$，所需的精确性是 $F(X)\equiv\exists Y\,A(X,Y)$：辅助状态可以隐藏，原始选择不能丢失。这个等式同时保证每个原始模型都能扩展到表示中，以及表示不会引入新的原始模型。第~\ref{sec:compilation}节将对具体编译表示证明这种双向保证。
